\section{Representing B Obligations in Lean}
\label{sec:encoding}

\subsection{B developments and the translation pipeline}

A B machine declares sets and constants, state variables, an invariant,
an initialization, and operations. Its proof obligations establish, among
other properties, that initialization satisfies the invariant and that
each operation preserves it under its precondition. Refinement introduces
more concrete state together with a gluing invariant and corresponding
proof obligations. \barrel accepts machines, refinements, implementations,
or an already generated POG file. It invokes Atelier~B's BXML and
proof-obligation generators when generation is needed.

The POG XML format~\cite{POGXML} records mathematical obligations, their
hypotheses, definitions, and typing information. \barrel parses this
format, extracts the context of each goal, and constructs a Lean
expression. For variables $x_1,\ldots,x_n$, hypotheses
$H_1,\ldots,H_m$, and conclusion $G$, the target has the schematic form
\[
  \forall x_1\cdots x_n,\quad
  \enc{H_1}\to\cdots\to\enc{H_m}\to\enc{G}.
\]
The translation also constructs the proof arguments required by partial
operators, described below. Unsupported expressions are reported as
encoding failures and prevent completion of the import.

\subsection{Sets, relations, and functions}

B's mathematical language uses sets, maplets, relational operators, and
quantified predicates. A relation $R$ from $A$ to $B$ is a subset of
$A\times B$. A partial function is a functional relation: whenever
$(x,y)\in R$ and $(x,z)\in R$, we have $y=z$. B writes
$R\in A\pfun B$ for membership in this partial-function space. Totality additionally
requires its domain to equal $A$. Thus a B function can also occur as a
relation in domain restriction, image, override, and other set operations.

\barrel uses Mathlib's \texttt{Set} and \texttt{SetRel} representations:
\[
  \mathsf{Set}\ \alpha=\alpha\to\mathsf{Prop},\qquad
  \mathsf{SetRel}\ \alpha\ \beta=\mathsf{Set}(\alpha\times\beta).
\]
B typing information determines the corresponding element types;
integers use Lean's \texttt{Int}, products use product types, and
powersets use \texttt{Set} over the element type. Function properties
are predicates on relations. This lets the same translated object be
used for both function application and relational reasoning.

In the current reader, unknown carrier-type identifiers are mapped to
integers, and deferred carrier sets receive hypotheses stating that
they are finite nonempty integer sets. The case study's sensor and zone
carriers are consequently represented by constrained integer sets,
rather than arbitrary Lean type parameters. The operator library also
provides finite sets, intervals, relational images, extrema, cardinality,
and application, with B-style notation and lemmas for interactive proof.
Operator coverage is an implementation property; acceptance of some
set-based constructs does not imply support for the entire B language.

\subsection{Partial operators with explicit evidence}

A B expression such as $\min(S)$ requires $S$ to have a least element.
Translating it into a total function on all Lean sets would leave a value
available even outside this intended domain. For extrema, application,
sequence size, and cardinality, \barrel instead includes a proof of the
precondition in the operator's type. For minimum on a partially ordered
carrier, the predicate and operator type are
\begin{align*}
  \WD_{\min}(S)&\;\equiv\;
    \exists x\in S,\ \forall y\in S,\ x\leq y,\\
  \mathsf{bmin}&:(S:\mathsf{Set}\ \alpha)\to
    \WD_{\min}(S)\to\alpha.
\end{align*}
The implementation uses the witnessed least element and proves membership
and minimality lemmas. On finite nonempty integer sets, the condition
follows from standard order properties. On an arbitrary ordered carrier,
nonemptiness alone is insufficient. Cardinality requires finiteness;
function application requires the evidence expressed by the library's
application predicate.

An operator occurrence therefore has an explicit proof dependency. During
translation, the encoder creates a fresh metavariable for a missing
proof argument and inserts it into the expression. It then collects the
remaining proof metavariables as WD obligations. Each condition is closed
over the variables and hypotheses available at that occurrence. Once the
condition is proved, the corresponding theorem supplies the expression's
proof argument. Users invoking the same library operator in subsequent
interactive proofs must also supply its evidence.

This discipline is not yet uniform across arithmetic operators. Division,
modulo, and exponentiation currently use direct Lean operators without
the same generated proof arguments. The paper's WD account concerns the
listed proof-carrying operators; the case study uses those
operators and ordinary integer comparisons. The trust boundary of the
translation is discussed in Section~\ref{sec:workflow}.

\subsection{Logical scope of WD evidence}

The position of a partial expression determines which assumptions may
justify it. For example, a reading $r(s)$ can occur after a guard
$s\in\operatorname{dom}(r)$, inside a quantifier whose bound variable is
$s$. Neither the guard nor the variable is a global assumption. The
encoder uses Lean's local context to retain precisely this scope while
constructing the expression and its missing proof arguments.

Write $\Gamma$ for the current context and $\enc{P}_\Gamma$ for a
predicate translated in that context. Implication and conjunction are
constructed as follows:
\begin{align*}
  \enc{P\Rightarrow Q}_\Gamma
    &=\forall h:\enc{P}_\Gamma,\ \enc{Q}_{\Gamma,h},\\
  \enc{P\land Q}_\Gamma
    &=\mathsf{DepAnd}\bigl(\lambda h:\enc{P}_\Gamma.\,
      \enc{Q}_{\Gamma,h}\bigr).
\end{align*}
Here $\mathsf{DepAnd}(Q)$ is defined as $\exists h:P,\ Q(h)$, for
$Q:P\to\mathsf{Prop}$. If $Q$ does not depend on $h$, this is equivalent
to ordinary conjunction. The dependent form allows the right conjunct
to contain operators whose evidence uses the left conjunct. The encoder
first translates $P$, then introduces $h$, and only then translates $Q$.
Consequently, a WD condition created within $P$ cannot use $P$ itself;
one created within $Q$ can. This order matters to the generated
conditions, even when reordering the predicates would preserve their
truth whenever all expressions are defined.
Before encoding, normalization splits conjunctive antecedents into
successive implications, preserving their order and exposing their
components as individual assumptions.

For both universal and existential quantification, the encoder introduces
the bound variable with its translated element type before visiting the
body. B set-membership restrictions remain predicates: introducing an
integer variable alone does not establish membership in a sensor set.
A condition created in the body is closed universally over its local
variables and proof assumptions, regardless of whether the source
quantifier is universal or existential. Thus an existential expression
does not by itself supply a particular witness for proving definedness.
For example, under the outer assumption $r\in S\pfun D$, the application in
\[
  \exists s,\quad s\in\operatorname{dom}(r)\land r(s)\le c
\]
requires a condition of the form
\[
  \forall s,\quad s\in\operatorname{dom}(r)
  \to\WD_{\mathrm{app}}(r,s),
\]
with the outer variables and assumptions also quantified. It does not
require that a suitable reading below $c$ exists: that remains the
principal existential goal. Disjunction and biconditional are encoded
with ordinary Lean connectives and introduce no corresponding local
hypothesis for either operand.

\subsection{A contextual WD obligation}

Consider a partial reading map $r$ and a set $A$ of available sensors
belonging to one zone. A lower-bound check uses $\min(r[A])$, where
$r[A]$ is the relational image of $A$. Its WD obligation must establish
that this image has a least element. If the surrounding context contains
\[
  r\in\mathit{SENSOR}\pfun\mathbb Z,
  \qquad A\in\operatorname{FIN}_1(\operatorname{dom}(r)),
\]
then the image is finite and nonempty. Here $\operatorname{FIN}_1(X)$ is
the collection of finite nonempty subsets of $X$. A library lemma proves
\[
  \bigl(r\in\mathit{SENSOR}\pfun\mathbb Z\bigr)
  \to\bigl(A\in\operatorname{FIN}_1(\operatorname{dom}(r))\bigr)
  \to\WD_{\min}(r[A]).
\]
The encoder's actual closed goal additionally quantifies the variables
and retains the hypotheses of the source occurrence. The proof may
establish the finite-nonempty premise from these hypotheses before
applying the lemma.

The context is essential. A summary operation guarded by $A\ne\varnothing$
can justify the minimum; an unguarded summary on empty readings cannot.
The encoding of implication and dependent conjunction allows subsequent
expressions to use proofs of preceding assumptions. The generated WD
statement records those dependencies, rather than merely asking whether
the syntactic expression $\min(r[A])$ is meaningful in isolation.

The resulting WD theorem also does not prove the principal property.
For example, proving that $\min(r[A])$ exists does not show that it exceeds
a configured threshold. The former belongs to definedness reasoning;
the latter is part of the authorization condition and its preservation
proof. This distinction structures both the automation and the case study.

\paragraph{From local evidence to named obligations.}
Let a missing argument have type $W$ in a local context
$x_1:T_1,\ldots,x_n:T_n$, where some $T_i$ are propositions. Its named
obligation has the closed type
\[
  d_W:\forall x_1:T_1,\ldots,\forall x_n:T_n,\ W.
\]
The generated expression retains a reference to this declaration as its
evidence, instantiated in the context of the occurrence. Nested partial
operators may also make one WD declaration depend on another. The importer
therefore instantiates all remaining proof metavariables before retaining
the generated statements. A principal goal mentioning an unproved WD
declaration remains pending: the individual proof command reports that
dependency and requires its proof first. The principal statement and the
WD statements thus form one proof development, rather than unrelated
lists of assertions.

This describes the implementation's evidence discipline for the supported
operators. Once every referenced declaration has a proof, Lean checks the
resulting terms against their types. It does not follow from that check
alone that the translation preserves B semantics; no mechanically verified
correctness theorem for the reader and encoder is claimed here.

\subsection{Sharing repeated conditions}

Different obligations can contain the same closed WD proposition.
\barrel maintains a table of conditions already encountered in the
current import, including both proved and pending ones. Canonicalization
removes binder names and expression metadata; a hash selects a bucket,
and comparison uses the canonical form with a bucket-local definitional
equality fallback. A match reuses the first condition's declaration
instead of creating another proof task.

The context remains part of the key. Conditions justified by different
hypotheses need not be shared even if they mention the same expression.
The procedure recognizes repeated representations, rather than proving
logical equivalence or subsumption. It does not share across imports,
and its hash-bucket restriction is not a complete test of definitional
equivalence. When a shared condition is still pending, its dependent
obligations continue to require its proof. Sharing removes duplicated
work without removing the common precondition.
